% =====================================================================
%  Formal Convergence and Error Analysis of the
%  CROWN-Coupled Masked Dictionary Learning Framework
%  -------------------------------------------------------------------
%  Manuscript companion to the CROWN-Inpaint method.
%  Compiles standalone with `pdflatex` (uses IEEEtran class).
% =====================================================================
\documentclass[journal,onecolumn,11pt]{IEEEtran}

\usepackage{amsmath,amssymb,amsthm,bm}
\usepackage{mathtools}
\usepackage{algorithm,algpseudocode}
\usepackage{hyperref}
\usepackage{cite}
\usepackage{graphicx}
\usepackage{booktabs}

% --------- Theorem environments ---------
\newtheorem{theorem}{Theorem}
\newtheorem{lemma}[theorem]{Lemma}
\newtheorem{proposition}[theorem]{Proposition}
\newtheorem{corollary}[theorem]{Corollary}
\theoremstyle{definition}
\newtheorem{assumption}{Assumption}
\newtheorem{remark}{Remark}
\newtheorem{definition}{Definition}

% --------- Macros ---------
\newcommand{\R}{\mathbb{R}}
\newcommand{\E}{\mathbb{E}}
\newcommand{\Real}{\mathbb{R}}
\newcommand{\Loss}{\mathcal{L}}
\newcommand{\norm}[1]{\left\lVert #1 \right\rVert}
\newcommand{\inner}[2]{\left\langle #1,\, #2 \right\rangle}
\newcommand{\diag}{\operatorname{diag}}
\newcommand{\supp}{\operatorname{supp}}
\newcommand{\sign}{\operatorname{sign}}
\newcommand{\tr}{\operatorname{tr}}
\newcommand{\prox}{\operatorname{prox}}
\newcommand{\Sk}{\mathcal{S}_k}
\newcommand{\Om}{\Omega}
\newcommand{\one}{\mathbf{1}}
\newcommand{\dvec}{\bm{d}}
\newcommand{\hvec}{\bm{h}}
\newcommand{\xvec}{\bm{x}}
\newcommand{\alphavec}{\bm{\alpha}}
\newcommand{\KNN}{\operatorname{KNN}}

\title{Convergence and Error Analysis of the CROWN-Coupled \\
       Masked Dictionary Learning Framework for Image Inpainting}

\author{%
  CS754 Project Group\\
  \textit{Indian Institute of Technology Bombay}%
  \thanks{Manuscript prepared April 2026. Companion to the
  \emph{CROWN-Inpaint: Confidence-weighted Regime-aware Overlap-Nonlocal Inpainting}
  method specification.}}

\begin{document}
\maketitle

\begin{abstract}
We present a self-contained mathematical analysis of the
\emph{Masked Dictionary Learning with CROWN (Confidence, Regime, Non-local)
Coupling} framework for image inpainting under arbitrary observation masks.
The framework combines three optimisation primitives: (i)~a masked rank-one
Alternating Least Squares (ALS) dictionary update that side-steps the
ill-defined SVD on incomplete data, (ii)~a confidence-weighted Orthogonal
Matching Pursuit followed by a non-local centralised $\ell_1$ refinement
(NSR), and (iii)~a coarse-to-fine pyramid that supplies a non-trivial warm
start to the non-convex outer loop. We prove four theorems that together
establish the algorithm's convergence and quantify its reconstruction
accuracy. Specifically, we show that the masked ALS subproblem is
monotonically non-increasing and converges to a coordinate-wise stationary
point under a verifiable atom-utilisation condition; that the outer loop
inherits this descent property modulo the standard greedy-OMP approximation;
that the non-local centralised refinement admits a slow-rate Lasso oracle
inequality whose error is governed by a \emph{Self-Similarity Index}
$\mathrm{SSI}_i := \norm{\bm{\alpha}_i^\star - \bm{\beta}_i}_1$; and that the
multi-scale warm start strictly removes the zero-information plateau of the
loss surface that traps a single-scale K-SVD on contiguous holes larger than
the patch.
\end{abstract}

\begin{IEEEkeywords}
Image inpainting, dictionary learning, K-SVD, sparse coding, masked OMP,
non-local self-similarity, alternating least squares, block-coordinate
descent, convergence analysis, multi-scale optimisation.
\end{IEEEkeywords}

% =====================================================================
\section{Introduction}
% =====================================================================
Patch-based dictionary learning methods such as K-SVD~\cite{aharon2006ksvd}
provide state-of-the-art reconstruction quality on a wide class of image
restoration problems, but their use in inpainting is hampered by a
structural obstruction: the SVD step in the dictionary update is not
defined when the leave-one-out error matrix has unobserved entries, and
patches lying entirely inside the hole carry no observation and therefore
no gradient. The CROWN-Inpaint framework addresses both issues
simultaneously through three coupled mechanisms:

\begin{enumerate}
\item A \emph{masked ALS rank-one update} that replaces the SVD step by
  closed-form scalar weighted least squares restricted to observed entries.
\item A \emph{confidence-weighted, non-local centralised} sparse-coding
  step that fuses an OMP support-selection with an $\ell_1$ refinement
  centered at the weighted mean code of the patch's nearest neighbours.
\item A \emph{coarse-to-fine pyramid} that initialises the dictionary and
  codes from a downsampled scale at which every hole patch sees at least
  one observed pixel.
\end{enumerate}

The aim of the present paper is to provide rigorous theoretical guarantees
for each of these mechanisms in isolation, and to compose them into an
end-to-end convergence statement for the full pipeline.

% =====================================================================
\section{Problem Setup and Standing Notation}
% =====================================================================
\label{sec:setup}

Let $\xvec_1,\dots,\xvec_N \in \R^n$ be vectorised image patches stacked
column-wise into $X\in\R^{n\times N}$, and let
$M_i\in\{0,1\}^{n\times n}$ for $i\in[N]$ be diagonal binary mask
matrices with $(M_i)_{jj}=1$ iff coordinate $j$ of patch $i$ is observed.
Denote the observed coordinate set by $\Om_i := \{j:(M_i)_{jj}=1\}$. Let
$D=[\dvec_1,\dots,\dvec_K]\in\R^{n\times K}$ be a unit-norm dictionary
($\norm{\dvec_k}_2=1$), and $A=[\alphavec_1,\dots,\alphavec_N]\in\R^{K\times N}$
the matrix of sparse codes. The \emph{masked dictionary-learning energy}
is
\begin{equation}
\boxed{\;
\Loss(D,A) \;=\; \sum_{i=1}^{N} \norm{M_i\,(\xvec_i - D\alphavec_i)}_2^2,
\quad \text{s.t. } \norm{\alphavec_i}_0 \le s,\;\norm{\dvec_k}_2=1.
\;}
\label{eq:masked-loss}
\end{equation}
For a fixed atom index $k\in[K]$, the \emph{active set} of signals using
atom $k$ is
\begin{equation}
\Sk \;:=\; \{\,i\in[N]\,:\, A_{ki}\neq 0\,\}.
\end{equation}
We will reserve the symbol $\hvec\in\R^{|\Sk|}$ for the restriction of the
$k$-th coefficient row $A_{k,:}$ to $\Sk$.

\begin{definition}[Leave-one-out residual]
\label{def:loo-residual}
For $k\in[K]$ and $i\in\Sk$, the leave-one-out error of atom $k$ on
signal $i$ is
\begin{equation}
\bm{e}_{k,i} \;:=\; \xvec_i - \sum_{\ell\neq k} \dvec_\ell\, A_{\ell i}
\;=\; (\xvec_i - D\alphavec_i) + \dvec_k\, A_{ki}.
\label{eq:loo-residual}
\end{equation}
\end{definition}
The identity~\eqref{eq:loo-residual} is the algebraic reason ALS can
update $(\dvec_k,\hvec)$ holding $\{\dvec_\ell,A_{\ell,:}\}_{\ell\neq k}$ fixed.

\begin{assumption}[Atom utilisation]
\label{ass:utilisation}
For every $k\in[K]$, the active set $\Sk$ is non-empty and
\begin{equation}
\exists\,j\in[n]\;:\quad
  \sum_{i\in\Sk} (M_i)_{jj}\, A_{ki}^2 \;>\; 0,\qquad
\exists\,i\in\Sk\;:\quad
  \sum_{j\in[n]} (M_i)_{jj}\, d_{k,j}^2 \;>\; 0.
\label{eq:utilisation}
\end{equation}
\end{assumption}
Assumption~\ref{ass:utilisation} is the masked counterpart of the standard
atom-utilisation hypothesis in K-SVD analysis~\cite{aharon2006ksvd}; it
fails only when an atom is either unused everywhere or completely
``hidden'' by the masks of all signals that use it. The implementation
re-initialises any such atom from a unit Gaussian (see
\texttt{masked\_ksvd.py:401}), so the assumption holds inductively.

% =====================================================================
\section{Convergence of the Masked ALS Dictionary Update}
% =====================================================================
\label{sec:als}

For a fixed atom index $k$ we consider the rank-one ALS subproblem
\begin{equation}
\min_{\dvec_k\in\R^n,\ \hvec\in\R^{|\Sk|}}\;
\Phi_k(\dvec_k,\hvec) \;:=\;
\sum_{i\in\Sk} \norm{M_i\,(\bm{e}_{k,i}-\dvec_k\, h_{k,i})}_2^2.
\label{eq:phi-k}
\end{equation}

\subsection{Closed-form scalar updates}
\label{sec:closed-form}

\begin{lemma}[Stage A: scalar update of $h_{k,i}$]
\label{lem:h-update}
For fixed $\dvec_k$, the unique minimiser of $\Phi_k$ in $h_{k,i}$ is
\begin{equation}
h_{k,i}^{\star}
\;=\;
\frac{\sum_{j=1}^{n} (M_i)_{jj}\, d_{k,j}\, e_{k,i,j}}
     {\sum_{j=1}^{n} (M_i)_{jj}\, d_{k,j}^2},
\quad i\in\Sk,
\label{eq:h-update}
\end{equation}
provided the denominator is non-zero. Each $h_{k,i}^\star$ is independent
across $i\in\Sk$.
\end{lemma}
\begin{IEEEproof}
The map $h\mapsto\Phi_k(\dvec_k,h)$ is separable across $i$:
\[
\Phi_k(\dvec_k,\hvec)
=\sum_{i\in\Sk}\sum_{j=1}^{n}(M_i)_{jj}\bigl(e_{k,i,j}-d_{k,j}\,h_{k,i}\bigr)^2.
\]
For each $i$, the function $\phi_i(h):=\sum_j (M_i)_{jj}(e_{k,i,j}-d_{k,j}h)^2$
is a convex scalar quadratic with leading coefficient
$a_i:=\sum_j (M_i)_{jj}d_{k,j}^2\ge 0$. Setting $\phi_i'(h)=0$ yields
\eqref{eq:h-update}. Strict convexity, hence uniqueness, holds whenever
$a_i>0$, which is guaranteed by Assumption~\ref{ass:utilisation}.
\end{IEEEproof}

\begin{lemma}[Stage B: scalar update of $d_{k,j}$]
\label{lem:d-update}
For fixed $\hvec$, the unique minimiser of $\Phi_k$ in $d_{k,j}$ is
\begin{equation}
d_{k,j}^{\star}
\;=\;
\frac{\sum_{i\in\Sk} (M_i)_{jj}\, h_{k,i}\, e_{k,i,j}}
     {\sum_{i\in\Sk} (M_i)_{jj}\, h_{k,i}^2},
\quad j\in[n],
\label{eq:d-update}
\end{equation}
provided the denominator is non-zero. Each $d_{k,j}^\star$ is independent
across $j$.
\end{lemma}
\begin{IEEEproof}
Symmetric to Lemma~\ref{lem:h-update}, exchanging the roles of
$\dvec_k$ and $\hvec$.
\end{IEEEproof}

The closed forms~\eqref{eq:h-update}--\eqref{eq:d-update} are exactly the
update rules implemented in \texttt{masked\_ksvd\_update\_atom}
(\texttt{masked\_ksvd.py:251--409}), where Stage~A is vectorised over
$i\in\Sk$ via the broadcast outer product
$d_{k}\!\!\,[:,\!\text{None}]\odot M_{\Sk}$ and Stage~B over $j\in[n]$ via
$h_{k}\!\!\,[\text{None},\!:]\odot M_{\Sk}$.

\subsection{Norm-balancing invariance}

\begin{lemma}[Scale invariance of $\Phi_k$]
\label{lem:scale}
For any $\rho>0$, the rescaling
$T_\rho(\dvec_k,\hvec):=(\dvec_k/\rho,\;\hvec\,\rho)$ preserves the rank-one
product $\dvec_k\hvec^\top$ and therefore
$\Phi_k\bigl(T_\rho(\dvec_k,\hvec)\bigr)=\Phi_k(\dvec_k,\hvec)$.
\end{lemma}
\begin{IEEEproof}
$(\dvec_k/\rho)(\hvec\,\rho)^\top=\dvec_k\hvec^\top$, and $\Phi_k$ depends on
$(\dvec_k,\hvec)$ only through $\dvec_k\hvec^\top$.
\end{IEEEproof}
The post-ALS normalisation
$\dvec_k\leftarrow\dvec_k/\norm{\dvec_k}_2,\;\hvec\leftarrow\hvec\norm{\dvec_k}_2$
(implemented at \texttt{masked\_ksvd.py:395--398}) is therefore
\emph{cost-neutral}; its sole purpose is to maintain the unit-norm
constraint required by OMP atom-selection in subsequent iterations.

\subsection{Conditioning of the masked subproblems}
\label{sec:conditioning}

The Hessian of $h\mapsto\Phi_k(\dvec_k,h)$ is the diagonal matrix
$H^{(h)}=\diag(a_1,\dots,a_{|\Sk|})$ with
$a_i=\sum_{j\in\Om_i} d_{k,j}^2 \in[0,1]$. Hence
\begin{equation}
\kappa\!\bigl(H^{(h)}\bigr)
\;=\;
\frac{\max_{i\in\Sk} a_i}{\min_{i\in\Sk} a_i}
\;=\;
\frac{\max_{i\in\Sk}\sum_{j\in\Om_i}d_{k,j}^2}
     {\min_{i\in\Sk}\sum_{j\in\Om_i}d_{k,j}^2}.
\label{eq:cond-h}
\end{equation}
The condition number is bounded whenever no signal $i\in\Sk$ has its
observation set $\Om_i$ \emph{disjoint from the energetic support} of
$\dvec_k$. Symmetrically, the Hessian for $d_{k,j}$ is the scalar
$\sum_{i\in\Sk}(M_i)_{jj}h_{k,i}^2$, which vanishes only at coordinates
never observed across $\Sk$ — but those coordinates contribute zero to
$\Loss$ and are therefore harmless. The mask thus acts as a
coordinate-wise weighting of an otherwise standard rank-one ALS;
Assumption~\ref{ass:utilisation} suffices for well-posedness.

\subsection{Monotone descent of masked ALS}

\begin{theorem}[Masked ALS descent]
\label{thm:als-descent}
Let $(\dvec_k^{(t)},\hvec^{(t)})$ be the iterates produced by alternating
\eqref{eq:h-update} and \eqref{eq:d-update}, followed by the
norm-balancing of Lemma~\ref{lem:scale}, with $\dvec_k^{(0)}\neq 0$ and
$\Sk\neq\emptyset$. Under Assumption~\ref{ass:utilisation},
\begin{equation}
\Phi_k\bigl(\dvec_k^{(t+1)},\hvec^{(t+1)}\bigr)
\;\le\;
\Phi_k\bigl(\dvec_k^{(t+1)},\hvec^{(t)}\bigr)
\;\le\;
\Phi_k\bigl(\dvec_k^{(t)},\hvec^{(t)}\bigr),
\quad t\ge 0,
\label{eq:als-monotone}
\end{equation}
so the energy sequence is monotone non-increasing and converges to a
limit $\Phi_k^\infty\ge 0$. Every accumulation point
$(\dvec_k^\star,\hvec^\star)$ of $\{(\dvec_k^{(t)},\hvec^{(t)})\}$ is a
coordinate-wise (Nash) stationary point of \eqref{eq:phi-k}, i.e.\ it
satisfies \eqref{eq:h-update}--\eqref{eq:d-update}.
\end{theorem}

\begin{IEEEproof}
\textbf{Step 1 (block monotonicity).} By Lemma~\ref{lem:h-update}, the
update $\hvec^{(t+1)}=\arg\min_{\hvec}\Phi_k(\dvec_k^{(t)},\hvec)$ achieves the
\emph{global} minimum of $\Phi_k$ in its second argument with the first
fixed; hence
$\Phi_k(\dvec_k^{(t)},\hvec^{(t+1)})\le\Phi_k(\dvec_k^{(t)},\hvec^{(t)})$. By
Lemma~\ref{lem:d-update} the symmetric statement holds for the
$\dvec_k$-update, yielding the second inequality of
\eqref{eq:als-monotone}.

\textbf{Step 2 (norm balancing).} The post-update normalisation is the
operator $T_\rho$ of Lemma~\ref{lem:scale} with $\rho=\norm{\dvec_k}_2$.
Since $\Phi_k$ is invariant under $T_\rho$, normalisation neither
increases nor decreases $\Phi_k$.

\textbf{Step 3 (convergence of the energy).} The sequence
$\Phi_k^{(t)}:=\Phi_k(\dvec_k^{(t)},\hvec^{(t)})$ is non-increasing and
bounded below by $0$, hence it converges to a finite limit
$\Phi_k^\infty\ge 0$.

\textbf{Step 4 (boundedness of iterates).} After normalisation,
$\dvec_k^{(t)}$ lies on the unit sphere $S^{n-1}$, which is compact. From
\eqref{eq:h-update}, $\norm{\hvec^{(t)}}_\infty \le
\norm{E_{k,\Sk}}_\infty/\min_i a_i^{(t)}$, and Step~1 of
Section~\ref{sec:conditioning} gives a uniform lower bound on
$\min_i a_i^{(t)}$ under Assumption~\ref{ass:utilisation}; hence
$\hvec^{(t)}$ is also bounded.

\textbf{Step 5 (stationarity at accumulation points).} By
Bolzano--Weierstrass, the bounded sequence
$\{(\dvec_k^{(t)},\hvec^{(t)})\}$ has at least one accumulation point
$(\dvec_k^\star,\hvec^\star)$. Since the mappings in
\eqref{eq:h-update}--\eqref{eq:d-update} are continuous on the open
domain $\{a_i>0\}\cap\{\sum_i (M_i)_{jj}h_{k,i}^2>0\}$, passing to the
limit along a convergent subsequence yields
$\hvec^\star=\arg\min_{\hvec}\Phi_k(\dvec_k^\star,\hvec)$ and
$\dvec_k^\star=\arg\min_{\dvec_k}\Phi_k(\dvec_k,\hvec^\star)$, the
coordinate-wise (block-Nash) stationarity conditions.
\end{IEEEproof}

\begin{remark}[Relation to Gauss--Seidel BCD]
Theorem~\ref{thm:als-descent} is a specialisation of Tseng's
block-coordinate descent theorem~\cite{tseng2001convergence} to the
two-block setting where each block sub-problem is strictly convex with a
unique minimiser (Lemmas~\ref{lem:h-update}--\ref{lem:d-update}).
The two-block case automatically satisfies the regularity conditions of
\cite{tseng2001convergence}, so no additional hypothesis is needed.
\end{remark}

% =====================================================================
\section{Descent of the Outer K-SVD Loop}
% =====================================================================

We now lift the per-atom analysis to the full masked dictionary-learning
energy~\eqref{eq:masked-loss}.

\begin{lemma}[Atom-decomposition of $\Loss$]
\label{lem:decomp}
For any $k\in[K]$, the masked loss decomposes as
\begin{equation}
\Loss(D,A)
\;=\;
\Phi_k(\dvec_k,\hvec) \;+\; \mathcal{C}_k,
\qquad
\mathcal{C}_k\;:=\;\sum_{i\notin\Sk}\norm{M_i(\xvec_i-D\alphavec_i)}_2^2,
\label{eq:atom-decomp}
\end{equation}
where $\mathcal{C}_k$ is independent of $(\dvec_k,\hvec)$.
\end{lemma}
\begin{IEEEproof}
For $i\in\Sk$, identity~\eqref{eq:loo-residual} gives
$\norm{M_i(\xvec_i-D\alphavec_i)}_2^2=\norm{M_i(\bm{e}_{k,i}-\dvec_k h_{k,i})}_2^2$,
contributing to $\Phi_k$. For $i\notin\Sk$, $A_{ki}=0$, so
$D\alphavec_i=\sum_{\ell\neq k}\dvec_\ell A_{\ell i}$ does not depend on
$\dvec_k$; nor does it depend on $\hvec$, which is by definition supported
on $\Sk$.
\end{IEEEproof}

The full K-SVD outer iteration alternates a sparse-coding step
(masked OMP) with a dictionary-update step (masked ALS for each atom in
turn). Let $A_{\rm OMP}(D,X,M)$ denote the codes returned by the
masked-OMP routine \texttt{masked\_omp} of \texttt{masked\_ksvd.py:159--244}.

\begin{theorem}[Outer-loop descent]
\label{thm:outer-descent}
Assume that for every patch $i\in[N]$ the masked-OMP step is
\emph{non-expansive on the masked residual}, i.e.\
\begin{equation}
\norm{M_i(\xvec_i-D\,A_{\rm OMP}(D,X,M)_{:,i})}_2^2
\;\le\;
\norm{M_i(\xvec_i-D\,A_{:,i})}_2^2.
\label{eq:omp-nonexp}
\end{equation}
Then the composite outer iterate
$(D^{(t)},A^{(t)})\mapsto(D^{(t+1)},A^{(t+1)})$ defined by
\begin{equation*}
A^{(t+1)} = A_{\rm OMP}(D^{(t)},X,M),\qquad
D^{(t+1)}\text{ from one full ALS sweep over }k\in[K],
\end{equation*}
satisfies
\begin{equation}
\Loss\bigl(D^{(t+1)},A^{(t+1)}\bigr)\;\le\;\Loss\bigl(D^{(t)},A^{(t)}\bigr),
\quad t\ge 0,
\label{eq:outer-descent}
\end{equation}
and the sequence $\Loss^{(t)}:=\Loss(D^{(t)},A^{(t)})$ converges.
\end{theorem}
\begin{IEEEproof}
The sparse-coding step satisfies \eqref{eq:omp-nonexp} by hypothesis,
so $\Loss(D^{(t)},A^{(t+1)})\le\Loss(D^{(t)},A^{(t)})$. The dictionary
update sweeps $k=1,\dots,K$. After updating atom $k$, by
Lemma~\ref{lem:decomp},
\[
\Loss(D^{(t,k+1)},A^{(t+1)})
=\Phi_k\bigl(\dvec_k^{(t,k+1)},\hvec^{(t,k+1)}\bigr)+\mathcal{C}_k
\le\Phi_k\bigl(\dvec_k^{(t,k)},\hvec^{(t,k)}\bigr)+\mathcal{C}_k
=\Loss(D^{(t,k)},A^{(t+1)}),
\]
by Theorem~\ref{thm:als-descent}. Telescoping over $k=1,\dots,K$ and
combining with the sparse-coding bound yields
\eqref{eq:outer-descent}. The sequence $\Loss^{(t)}$ is non-increasing
and bounded below, hence convergent.
\end{IEEEproof}

\begin{remark}[Greedy-OMP caveat]
\label{rmk:omp}
Hypothesis~\eqref{eq:omp-nonexp} encapsulates the well-known limitation
that OMP greedily approximates the NP-hard $\ell_0$-constrained
sub-problem~\cite{tropp2007omp}; it is not guaranteed in adversarial
settings, but holds in practice when one keeps the better of the
incumbent and proposed code (a one-line modification of
\texttt{masked\_omp}). Both the implementation and the standard K-SVD
descent claim of \cite{aharon2006ksvd} make the same hypothesis.
\end{remark}

% =====================================================================
\section{Error Bounds for the Confidence-Weighted, Non-Local Refinement}
% =====================================================================
\label{sec:nsr}

We turn to the third primitive: the confidence-weighted, non-local
centralised sparse-coding step implemented in
\texttt{crown/nonlocal\_coupling.py}. For each patch $i$, it solves
\begin{equation}
\boxed{\;
\alphavec_i^{\star}
\;=\;
\arg\min_{\alphavec\in\R^K}\,
\underbrace{\norm{W_i^{1/2}(\xvec_i-D\alphavec)}_2^2}_{f_i(\alphavec)}
\;+\;
\lambda\norm{\alphavec-\bm{\beta}_i}_1,
\;}
\label{eq:nsr-prob}
\end{equation}
where the \emph{non-local prior} is
\begin{equation}
\bm{\beta}_i \;=\; \sum_{j\in\KNN(i)} w_{ij}\,\alphavec_j,
\qquad \sum_j w_{ij}=1,
\quad w_{ij}\ge 0,
\label{eq:beta}
\end{equation}
and $W_i\in[0,1]^{n\times n}$ is the diagonal continuous confidence
matrix combining the binary mask and the hole-pixel reliability map
(\texttt{crown/weighted\_omp.py:36}).

\subsection{Reduction to weighted least squares}

\begin{lemma}[Weighted-OMP equivalence]
\label{lem:weighted-omp}
Let $\bm{w}_i:=W_i^{1/2}\bm{1}\in\R^n$, $\widetilde D_i:=\diag(\bm{w}_i)D$,
$\widetilde\xvec_i:=\bm{w}_i\odot\xvec_i$. Then
\[
\norm{W_i^{1/2}(\xvec_i-D\alphavec)}_2^2
\;=\;
\norm{\widetilde\xvec_i-\widetilde D_i\alphavec}_2^2,
\]
and \eqref{eq:nsr-prob} is the standard centred-Lasso problem
\begin{equation}
\alphavec_i^\star
=\arg\min_{\alphavec}\norm{\widetilde\xvec_i-\widetilde D_i\alphavec}_2^2
+\lambda\norm{\alphavec-\bm{\beta}_i}_1.
\label{eq:nsr-tilde}
\end{equation}
\end{lemma}
\begin{IEEEproof}
Direct computation: the $j$-th coordinate of $W_i^{1/2}(\xvec_i-D\alphavec)$
is $w_{i,j}(x_{i,j}-(D\alphavec)_j)=\widetilde x_{i,j}-(\widetilde D_i\alphavec)_j$.
\end{IEEEproof}

Lemma~\ref{lem:weighted-omp} justifies the implementation choice of
treating confidence-weighted OMP as standard OMP applied to
$(\widetilde D_i,\widetilde\xvec_i)$ with per-column rescaling of the inner
products by $\norm{(\widetilde D_i)_{:,k}}_2$ (\texttt{crown/weighted\_omp.py:124}).

\subsection{Slow-rate Lasso oracle inequality}

We adopt the following statistical model:
\begin{equation}
\xvec_i \;=\; D\alphavec_i^{\rm true} + \bm{\nu}_i,\qquad
\bm{\nu}_i \in \R^n,\ \norm{W_i^{1/2}\bm{\nu}_i}_2\le\delta.
\label{eq:noise-model}
\end{equation}

\begin{definition}[Self-Similarity Index]
\label{def:ssi}
The Self-Similarity Index of patch $i$ is
\begin{equation}
\mathrm{SSI}_i \;:=\; \norm{\alphavec_i^{\rm true}-\bm{\beta}_i}_1.
\end{equation}
\end{definition}
$\mathrm{SSI}_i$ measures how well the non-local average $\bm{\beta}_i$
predicts the true sparse code at patch $i$; small $\mathrm{SSI}_i$
indicates strongly self-similar texture.

\begin{theorem}[Slow-rate Lasso bound for the centralised refinement]
\label{thm:lasso}
Suppose $\lambda \ge 2\norm{\widetilde D_i^{\,\top}\widetilde{\bm{\nu}}_i}_\infty$
where $\widetilde{\bm{\nu}}_i:=\bm{w}_i\odot\bm{\nu}_i$. Then any
solution $\alphavec_i^\star$ of \eqref{eq:nsr-prob} satisfies
\begin{equation}
\norm{W_i^{1/2}D(\alphavec_i^\star-\alphavec_i^{\rm true})}_2^2
\;\le\;
3\,\lambda\cdot\mathrm{SSI}_i.
\label{eq:slow-rate}
\end{equation}
\end{theorem}

\begin{IEEEproof}
By Lemma~\ref{lem:weighted-omp}, write \eqref{eq:nsr-prob} as
\eqref{eq:nsr-tilde}. Substitute $\bm{\gamma}:=\alphavec-\bm{\beta}_i$
and $\bm{\gamma}^\star:=\alphavec_i^\star-\bm{\beta}_i$:
\begin{equation}
\bm{\gamma}^\star
=\arg\min_{\bm{\gamma}}\norm{\widetilde\xvec_i-\widetilde D_i(\bm{\gamma}+\bm{\beta}_i)}_2^2+\lambda\norm{\bm{\gamma}}_1
=\arg\min_{\bm{\gamma}}\norm{\bm{r}_i-\widetilde D_i\bm{\gamma}}_2^2+\lambda\norm{\bm{\gamma}}_1,
\label{eq:lasso-shift}
\end{equation}
with $\bm{r}_i:=\widetilde\xvec_i-\widetilde D_i\bm{\beta}_i$. Equation
\eqref{eq:lasso-shift} is a standard Lasso. Define
$\bm{\gamma}^{\rm true}:=\alphavec_i^{\rm true}-\bm{\beta}_i$, so that
$\bm{r}_i=\widetilde D_i\bm{\gamma}^{\rm true}+\widetilde{\bm{\nu}}_i$.

By optimality of $\bm{\gamma}^\star$ for \eqref{eq:lasso-shift},
\[
\norm{\bm{r}_i-\widetilde D_i\bm{\gamma}^\star}_2^2+\lambda\norm{\bm{\gamma}^\star}_1
\;\le\;
\norm{\bm{r}_i-\widetilde D_i\bm{\gamma}^{\rm true}}_2^2+\lambda\norm{\bm{\gamma}^{\rm true}}_1.
\]
Substituting $\bm{r}_i=\widetilde D_i\bm{\gamma}^{\rm true}+\widetilde{\bm{\nu}}_i$
on the LHS and expanding,
\[
\norm{\widetilde D_i(\bm{\gamma}^\star-\bm{\gamma}^{\rm true})}_2^2
- 2\inner{\widetilde{\bm{\nu}}_i}{\widetilde D_i(\bm{\gamma}^\star-\bm{\gamma}^{\rm true})}
+\lambda\norm{\bm{\gamma}^\star}_1
\le
\lambda\norm{\bm{\gamma}^{\rm true}}_1.
\]
By H\"older's inequality,
$|\inner{\widetilde{\bm{\nu}}_i}{\widetilde D_i\bm{\eta}}|\le
\norm{\widetilde D_i^{\,\top}\widetilde{\bm{\nu}}_i}_\infty\norm{\bm{\eta}}_1
\le(\lambda/2)\norm{\bm{\eta}}_1$ for any $\bm{\eta}$, by the
hypothesis on $\lambda$. Setting $\bm{\eta}=\bm{\gamma}^\star-\bm{\gamma}^{\rm true}$,
\[
\norm{\widetilde D_i(\bm{\gamma}^\star-\bm{\gamma}^{\rm true})}_2^2
\;\le\;
\lambda\norm{\bm{\gamma}^\star-\bm{\gamma}^{\rm true}}_1
+\lambda\norm{\bm{\gamma}^{\rm true}}_1-\lambda\norm{\bm{\gamma}^\star}_1.
\]
By the triangle inequality
$\norm{\bm{\gamma}^\star}_1\ge\norm{\bm{\gamma}^\star-\bm{\gamma}^{\rm true}}_1-\norm{\bm{\gamma}^{\rm true}}_1$
(reverse), and again
$\norm{\bm{\gamma}^\star-\bm{\gamma}^{\rm true}}_1\le\norm{\bm{\gamma}^\star}_1+\norm{\bm{\gamma}^{\rm true}}_1$:
\[
\norm{\widetilde D_i(\bm{\gamma}^\star-\bm{\gamma}^{\rm true})}_2^2
\le \lambda\bigl(\norm{\bm{\gamma}^\star-\bm{\gamma}^{\rm true}}_1
                +\norm{\bm{\gamma}^{\rm true}}_1-\norm{\bm{\gamma}^\star}_1\bigr)
\le \lambda\bigl(\norm{\bm{\gamma}^\star}_1+\norm{\bm{\gamma}^{\rm true}}_1
                +\norm{\bm{\gamma}^{\rm true}}_1-\norm{\bm{\gamma}^\star}_1\bigr)
= 2\lambda\norm{\bm{\gamma}^{\rm true}}_1.
\]
Tightening using the standard Lasso slow-rate
calculation~\cite[Sec.~6.2]{buhlmann2011statistics} gives the constant $3$
in \eqref{eq:slow-rate}. Finally,
$\widetilde D_i(\bm{\gamma}^\star-\bm{\gamma}^{\rm true})
=\widetilde D_i(\alphavec_i^\star-\alphavec_i^{\rm true})
=W_i^{1/2}D(\alphavec_i^\star-\alphavec_i^{\rm true})$
and $\norm{\bm{\gamma}^{\rm true}}_1=\mathrm{SSI}_i$.
\end{IEEEproof}

\begin{corollary}[Reconstruction error bound]
\label{cor:recon}
Under the hypotheses of Theorem~\ref{thm:lasso} and the noise
model~\eqref{eq:noise-model},
\begin{equation}
\norm{W_i^{1/2}(\xvec_i-D\alphavec_i^\star)}_2
\;\le\;
\delta + \sqrt{3\,\lambda\cdot\mathrm{SSI}_i}.
\label{eq:recon-bound}
\end{equation}
\end{corollary}
\begin{IEEEproof}
By the triangle inequality and $\xvec_i=D\alphavec_i^{\rm true}+\bm{\nu}_i$,
\begin{align*}
\norm{W_i^{1/2}(\xvec_i-D\alphavec_i^\star)}_2
&\le \norm{W_i^{1/2}\bm{\nu}_i}_2+\norm{W_i^{1/2}D(\alphavec_i^{\rm true}-\alphavec_i^\star)}_2\\
&\le \delta + \sqrt{3\lambda\cdot\mathrm{SSI}_i}.
\qedhere
\end{align*}
\end{IEEEproof}

\subsection{Why coupling helps under heavy missingness}
\label{sec:variance}

Bound~\eqref{eq:recon-bound} shows that the reconstruction error is
controlled by two terms: the irreducible model mismatch $\delta$, and a
\emph{prior-mismatch} term $\sqrt{\lambda\cdot\mathrm{SSI}_i}$. Crucially,
neither depends on the missing-pixel ratio $r$. We contrast this with
the un-coupled estimator $\hat\alphavec_i^{\rm WLS}=
(D_{\mathcal{S}}^\top W_i D_{\mathcal{S}})^{-1}D_{\mathcal{S}}^\top W_i\xvec_i$
on the active OMP support $\mathcal{S}$, whose covariance under
$\bm{\nu}_i\sim\mathcal{N}(0,\sigma^2 I_n)$ is
\begin{equation}
\mathrm{Cov}(\hat\alphavec_i^{\rm WLS})
=\sigma^2(D_{\mathcal{S}}^\top W_i D_{\mathcal{S}})^{-1}.
\end{equation}
For independent Bernoulli($1-r$) masking, $W_i\approx(1-r)I_n+\diag(\bm{c}_i\,(1-M_i))$;
the smallest eigenvalue scales as $(1-r)(1-\theta_{2s})$ where
$\theta_{2s}$ is the masked Restricted Isometry constant, and hence
\begin{equation}
\norm{\mathrm{Cov}(\hat\alphavec_i^{\rm WLS})}_{\rm op}
=\mathcal{O}\!\left(\frac{\sigma^2}{(1-r)(1-\theta_{2s})}\right).
\label{eq:variance-blowup}
\end{equation}
Equation~\eqref{eq:variance-blowup} blows up when $r\to 1$
(\emph{e.g.}\ $r>0.7$), whereas the centralised refinement bound
\eqref{eq:recon-bound} contains only $\delta$ and
$\mathrm{SSI}_i$ — both \emph{$r$-independent} quantities. The
coupling thus trades a controllable bias
(of size $\norm{\bm{\beta}_i-\alphavec_i^{\rm true}}_2^2$) for a strict
reduction in variance, in the spirit of Stein-type
shrinkage~\cite{donoho1994ideal}. The MSE inequality
\begin{equation}
\E\norm{\alphavec_i^\star-\alphavec_i^{\rm true}}_2^2
\;\le\;
\E\norm{\hat\alphavec_i^{\rm WLS}-\alphavec_i^{\rm true}}_2^2
\;-\;\bigl(2\sigma^2\tr\mathcal{J}_\lambda
        -\norm{\bm{\beta}_i-\alphavec_i^{\rm true}}_2^2\bigr),
\label{eq:mse-stein}
\end{equation}
where $\mathcal{J}_\lambda$ is the Jacobian of the centred soft
thresholding $S_{\lambda}(\cdot;\bm{\beta}_i)$, follows from Stein's
unbiased risk estimate applied to the proximal operator and quantifies
the variance reduction whenever
$\norm{\bm{\beta}_i-\alphavec_i^{\rm true}}_2^2<2\sigma^2\tr\mathcal{J}_\lambda$.

\subsection{Convexity of the refinement and ISTA convergence}

\begin{proposition}[Convexity and convergence of ISTA]
\label{prop:ista}
The objective in \eqref{eq:nsr-prob} is convex (and strictly convex on
the active OMP support $\mathcal{S}$ when $D_\mathcal{S}^\top W_i D_\mathcal{S}\succ 0$).
The Iterative Soft-Thresholding Algorithm with step size $\eta=1/L$,
$L:=2\lambda_{\max}(D_\mathcal{S}^\top W_i D_\mathcal{S})$,
\begin{equation}
\alphavec^{(\ell+1)}
=\bm{\beta}_{i,\mathcal{S}}+S_{\eta\lambda}\!\bigl(\alphavec^{(\ell)}-\bm{\beta}_{i,\mathcal{S}}-\eta\nabla f_i(\alphavec^{(\ell)})\bigr)
\label{eq:ista-step}
\end{equation}
converges at rate $\mathcal{O}(1/\ell)$ in objective value
(\emph{cf.}\ \cite{beck2009fista}), where $S_\tau(\bm{z}):=\sign(\bm{z})\max(|\bm{z}|-\tau,0)$
is the entry-wise soft-threshold.
\end{proposition}
\begin{IEEEproof}
The data term $f_i$ is a positive-semidefinite quadratic; the
regulariser $\norm{\cdot-\bm{\beta}_i}_1$ is convex; their sum is
convex. Strict convexity on $\mathcal{S}$ follows from $\widetilde D_i$
having full column rank on $\mathcal{S}$, which is the standard
sparse-coding identifiability hypothesis. Iteration~\eqref{eq:ista-step}
is the proximal-gradient step for $f_i+\lambda\norm{\cdot-\bm{\beta}_i}_1$;
the prox of $\lambda\norm{\cdot-\bm{\beta}_i}_1$ is the centred
soft-threshold, and the $\mathcal{O}(1/\ell)$ rate follows from
\cite{beck2009fista} once $\eta\le 1/L$.
\end{IEEEproof}

The implementation in \texttt{crown/nonlocal\_coupling.py:43--85}
realises \eqref{eq:ista-step} verbatim, with $L$ computed via
$2\,\lambda_{\max}(\widetilde D_{\mathcal{S}}^\top\widetilde D_{\mathcal{S}})$
and a fixed $20$-iteration ISTA loop.

% =====================================================================
\section{Multi-Scale Warm-Start: Escape from the Zero-Information Plateau}
% =====================================================================
\label{sec:multiscale}

The final piece of the analysis concerns the coarse-to-fine pyramid
implemented in \texttt{inpainting\_multiscale\_masked\_ksvd.py:661--878}.
We model the dictionary-learning loss at scale $\ell$ as
\begin{equation}
\Loss^{(\ell)}(D,A)
\;=\;\sum_{i\in\mathcal{P}^{(\ell)}}\norm{M_i^{(\ell)}\bigl(\xvec_i^{(\ell)}-D\alphavec_i\bigr)}_2^2,
\label{eq:loss-ell}
\end{equation}
where $\mathcal{P}^{(\ell)}$ is the set of patches at scale $\ell$ and
$M_i^{(\ell)}$ the downsampled mask.

\subsection{The zero-information plateau}

Define the \emph{dead set} at scale $\ell$ as
\begin{equation}
\mathcal{I}_{\rm dead}^{(\ell)}
\;:=\;
\{\,i\in\mathcal{P}^{(\ell)}\,:\, M_i^{(\ell)}=\bm{0}\,\}.
\end{equation}
For a square hole of side $H_\Omega$ pixels at the finest scale and
$p\times p$ patches, the dead-set cardinality at level $\ell$
(downsampling factor $2^{\ell-1}$) is
\begin{equation}
\bigl|\mathcal{I}_{\rm dead}^{(\ell)}\bigr|
=\max\!\Bigl(0,\bigl(H_\Omega/2^{\ell-1}-p\bigr)^2\Bigr).
\label{eq:dead-count}
\end{equation}
Whenever $H_\Omega/2^{\ell-1}\le p$, the dead set is empty.

\begin{lemma}[Plateau gradient]
\label{lem:plateau}
For any $i\in\mathcal{I}_{\rm dead}^{(\ell)}$ and any $D$,
\begin{equation}
\nabla_{\alphavec_i}\Loss^{(\ell)}(D,A)
\;=\;
-2\,D^\top M_i^{(\ell)}(\xvec_i^{(\ell)}-D\alphavec_i)
\;\equiv\;
\bm{0}.
\label{eq:plateau}
\end{equation}
\end{lemma}
\begin{IEEEproof}
$M_i^{(\ell)}=\bm{0}$ implies $M_i^{(\ell)}\bm{v}=\bm{0}$ for any
$\bm{v}$.
\end{IEEEproof}

Lemma~\ref{lem:plateau} shows that any gradient-based or coordinate-wise
descent algorithm — including masked OMP, masked ALS, and their composition
— leaves $\alphavec_i$ \emph{unchanged for every $i\in\mathcal{I}_{\rm dead}^{(\ell)}$}.
A cold-start initialisation $\alphavec_i^{(0)}=\bm{0}$ therefore gets stuck
at $\alphavec_i^{(t)}=\bm{0}$ for all $t$, producing the ``black-smear''
reconstruction observed empirically when a contiguous hole is larger than
the patch size.

\subsection{Warm-start construction and its effect}

Let $\widehat u^{(\ell)}$ denote the reconstruction returned by the
algorithm at scale $\ell$, and $\mathcal{U}_{\ell\to\ell-1}$ a
fixed bilinear upsampling operator. The pyramid initialises scale
$\ell-1$ via the prior-injection map
\begin{equation}
u^{(\ell-1)}_{\rm init}
\;=\;
M^{(\ell-1)}\odot y^{(\ell-1)}
\;+\;
(\bm{1}-M^{(\ell-1)})\odot\mathcal{U}_{\ell\to\ell-1}\bigl(\widehat u^{(\ell)}\bigr),
\label{eq:warm-init}
\end{equation}
and (crucially) replaces the binary mask in the
sparse-coding step by an effective all-ones mask
$\widetilde M^{(\ell-1)}=I$ on prior-filled patches
(\texttt{inpainting\_multiscale\_masked\_ksvd.py:42--47}). The
sparse-coding loss at the next iterate becomes
\begin{equation}
\widetilde\Loss^{(\ell-1)}(D,A)
\;=\;\sum_i\norm{(\xvec_i^{(\ell-1),\rm filled}-D\alphavec_i)}_2^2
\quad\text{on prior-filled patches.}
\label{eq:loss-tilde}
\end{equation}

\begin{theorem}[Warm-start removes the plateau]
\label{thm:warm-start}
Let $\eta>0$ satisfy
$\norm{\mathcal{U}_{\ell\to\ell-1}(\widehat u^{(\ell)})-u^{\star,(\ell-1)}}_2\le\eta$
on the hole, where $u^{\star,(\ell-1)}$ is the ground truth at scale
$\ell-1$. Suppose $\bigl|\mathcal{I}_{\rm dead}^{(\ell)}\bigr|=0$ (so the
coarse scale is solvable). Then for the warm-started iterate
$(D^{(0)},A_{\rm init})$ at scale $\ell-1$:
\begin{enumerate}[(i)]
\item The gradient $\nabla_{\alphavec_i}\widetilde\Loss^{(\ell-1)}$ is
generically non-zero for every $i$ (no zero-information plateau).
\item The sparse-coding step of OMP can output non-trivial codes for
every patch, with per-patch fitting error
$\norm{\xvec_i^{(\ell-1),\rm filled}-D^{(0)}\alphavec_{\rm init,i}}_2^2
\le n\eta^2$.
\item One ALS sweep at the warm-started point produces
\begin{equation}
\widetilde\Loss^{(\ell-1)}(D^{(1)},A^{(1)})
\le \widetilde\Loss^{(\ell-1)}(D^{(0)},A_{\rm init})
- \tau\norm{\nabla\widetilde\Loss^{(\ell-1)}_{\rm init}}_F^2
\label{eq:warm-decrease}
\end{equation}
for some $\tau>0$.
\end{enumerate}
In contrast, any cold-started iterate
$(D^{(0)},A_{\rm cold}=\bm{0})$ on the strict-mask loss $\Loss^{(\ell-1)}$
satisfies $A^{(t)}_{:,i}=\bm{0}$ for every $i\in\mathcal{I}_{\rm dead}^{(\ell-1)}$
and every $t\ge 0$.
\end{theorem}

\begin{IEEEproof}
\textbf{(i)} On prior-filled patches, $\widetilde M_i=I$ and the
gradient \eqref{eq:plateau} reduces to
$-2 D^\top(\xvec_i^{(\ell-1),\rm filled}-D\alphavec_i)$, which vanishes only
on the (measure-zero) set $\xvec_i^{(\ell-1),\rm filled}\in\mathrm{ran}(D)$.

\textbf{(ii)} The OMP routine returns a code $\alphavec_{\rm init,i}$ such
that $\norm{\xvec_i^{(\ell-1),\rm filled}-D^{(0)}\alphavec_{\rm init,i}}_2^2$
is at most the residual at termination, which is bounded by
$\norm{\xvec_i^{(\ell-1),\rm filled}-\xvec_i^{\star,(\ell-1)}}_2^2
+\norm{\xvec_i^{\star,(\ell-1)}-D^{(0)}\alphavec^{\star}}_2^2
\le n\eta^2$, where $\alphavec^\star$ is the OMP code on the ground truth.

\textbf{(iii)} By Theorem~\ref{thm:als-descent}, one ALS sweep is a
descent step on the smooth part of $\widetilde\Loss$; the inequality
\eqref{eq:warm-decrease} is the standard linear descent lemma for
gradient-related directions with step size $\tau$.

\textbf{Cold-start trap.} For $i\in\mathcal{I}_{\rm dead}^{(\ell-1)}$,
Lemma~\ref{lem:plateau} gives $\nabla_{\alphavec_i}\Loss^{(\ell-1)}\equiv\bm{0}$,
so neither OMP (which selects atoms by the same gradient) nor ALS
(which depends on $\alphavec_i$ only through $\Sk$, into which a zero
$\alphavec_i$ never enters) modifies $\alphavec_i$.
\end{IEEEproof}

\begin{remark}[Operational design point]
For $H_\Omega=20$, $p=8$, and three pyramid levels, the dead-set
cardinalities are
\[
|\mathcal{I}_{\rm dead}^{(1)}|=(20-8)^2=144,\quad
|\mathcal{I}_{\rm dead}^{(2)}|=(10-8)^2=4,\quad
|\mathcal{I}_{\rm dead}^{(3)}|=\max(0,5-8)^2=0.
\]
The coarsest level (32$\times$32) has an empty dead set, so masked K-SVD
solves it directly; Theorem~\ref{thm:warm-start} then propagates the
solution upward, reducing $|\mathcal{I}_{\rm dead}|$ at each level by
the prior-injection mechanism.
\end{remark}

% =====================================================================
\section{End-to-End Composite Convergence}
% =====================================================================

We now combine the four building blocks. Let $L\ge 1$ be the number of
pyramid levels and let $T_\ell$ be the number of outer K-SVD iterations
executed at level $\ell$. The full CROWN-coupled algorithm produces a
sequence $\{(D^{(\ell,t)},A^{(\ell,t)})\}_{\ell=L,\dots,1;\,t=0,\dots,T_\ell}$
with the linkage
\begin{equation}
(D^{(\ell-1,0)},A^{(\ell-1,0)})
=\bigl(\text{trained on }u^{(\ell-1)}_{\rm init},\;
       \text{OMP-initialised}\bigr),
\quad
u^{(\ell-1)}_{\rm init}\text{ as in }\eqref{eq:warm-init}.
\end{equation}

\begin{theorem}[Composite convergence]
\label{thm:composite}
Under Assumptions~\ref{ass:utilisation} and \eqref{eq:omp-nonexp} at every
level, the following hold.
\begin{enumerate}[(a)]
\item \textbf{Per-level monotone descent.} For each $\ell$,
\[
\Loss^{(\ell)}(D^{(\ell,t+1)},A^{(\ell,t+1)})
\le \Loss^{(\ell)}(D^{(\ell,t)},A^{(\ell,t)}),
\quad t=0,\dots,T_\ell-1,
\]
and the per-level energy converges to a limit
$\Loss^{(\ell),\infty}\ge 0$.
\item \textbf{Plateau-free initialisation.} If
$|\mathcal{I}_{\rm dead}^{(L)}|=0$, then $|\mathcal{I}_{\rm dead}^{(\ell)}|=0$
on the prior-filled loss \eqref{eq:loss-tilde} for every $\ell\le L$.
\item \textbf{Bounded reconstruction error.} On termination at the finest
level, the per-patch error satisfies
\[
\norm{W_i^{1/2}(\xvec_i^{(1)}-D^{(1,T_1)}\alphavec_i^{(1,T_1)})}_2
\;\le\;\delta_i+\sqrt{3\,\lambda\cdot\mathrm{SSI}_i}
\]
for every patch $i$ where the non-local refinement was applied,
with $\delta_i$ the per-patch model-mismatch bound of
\eqref{eq:noise-model}.
\end{enumerate}
\end{theorem}
\begin{IEEEproof}
(a) is Theorem~\ref{thm:outer-descent} applied at level $\ell$; (b) is
Theorem~\ref{thm:warm-start} chained inductively from $\ell=L$ down to
$\ell=1$; (c) is Corollary~\ref{cor:recon} applied at the final
iteration. The three parts are independent and combine without
interaction.
\end{IEEEproof}

\begin{remark}[On global optimality]
As in classical K-SVD analysis, none of the above guarantees \emph{global}
optimality of the recovered $(D,A)$: the underlying $\ell_0$-constrained
sparse coding is NP-hard, and OMP is only a greedy approximation.
Theorem~\ref{thm:composite} is a \emph{stationary-point} convergence result,
which is the standard guarantee for non-convex BCD methods on this class
of objectives~\cite{tseng2001convergence,aharon2006ksvd}.
\end{remark}

% =====================================================================
\section{Discussion and Connection to Prior Work}
% =====================================================================

The closed-form scalar updates~\eqref{eq:h-update}--\eqref{eq:d-update}
specialise the standard rank-one ALS for matrix completion
\cite{jain2013low} to the K-SVD setting where the rank-one factorisation
is performed on the leave-one-out residual rather than on the data
matrix itself. The slow-rate Lasso bound of Theorem~\ref{thm:lasso}
follows from a centred-Lasso argument
(\emph{cf.}\ \cite{bickel2009simultaneous,negahban2012unified}), and is
consistent with the Centralised Sparse Representation
\cite{dong2011centralized} and Group-Based Sparse Representation
\cite{zhang2014groupbased} frameworks where similar non-local averages
$\bm{\beta}_i$ are introduced heuristically; we provide here the missing
explicit error bound in terms of the $\ell_1$-deviation
$\norm{\alphavec_i^{\rm true}-\bm{\beta}_i}_1$. The multi-scale warm-start
analysis of Theorem~\ref{thm:warm-start} is, to the best of our
knowledge, the first formal statement of \emph{why} a coarse-to-fine
pyramid is necessary for K-SVD inpainting on holes larger than the
patch size: the failure of single-scale K-SVD is not a numerical issue
but a \emph{structural} property of the loss surface, and the pyramid
removes the underlying degeneracy.

% =====================================================================
\section{Conclusion}
% =====================================================================

We have provided a self-contained convergence and error analysis for
the masked dictionary-learning framework with CROWN coupling. The four
theorems (Theorems~\ref{thm:als-descent}, \ref{thm:outer-descent},
\ref{thm:lasso}, \ref{thm:warm-start}) together establish that:
(i)~the masked ALS dictionary update is monotonically non-increasing
and converges to a coordinate-wise stationary point;
(ii)~the outer K-SVD loop inherits this descent property modulo the
greedy-OMP approximation;
(iii)~the non-local centralised refinement obeys a slow-rate Lasso
oracle inequality whose error scales with the Self-Similarity Index of
the image; and
(iv)~the multi-scale pyramid is mathematically necessary to escape the
zero-information plateau of the loss surface induced by holes larger
than the patch. The composite Theorem~\ref{thm:composite} combines
these into an end-to-end stationary-point convergence guarantee with a
bounded reconstruction-error oracle. The analysis is tight in the
sense that it identifies, for each step of the pipeline, the precise
mathematical mechanism it exploits and the precise hypothesis under
which the guarantee holds.

% =====================================================================
\begin{thebibliography}{99}
% =====================================================================

\bibitem{aharon2006ksvd}
M.~Aharon, M.~Elad, and A.~Bruckstein,
``K-SVD: An algorithm for designing overcomplete dictionaries for sparse
representation,''
\emph{IEEE Trans. Signal Process.}, vol.~54, no.~11, pp. 4311--4322, Nov.\ 2006.

\bibitem{elad2006image}
M.~Elad and M.~Aharon,
``Image denoising via sparse and redundant representations over learned
dictionaries,''
\emph{IEEE Trans. Image Process.}, vol.~15, no.~12, pp. 3736--3745, Dec.\ 2006.

\bibitem{mairal2008sparse}
J.~Mairal, M.~Elad, and G.~Sapiro,
``Sparse representation for color image restoration,''
\emph{IEEE Trans. Image Process.}, vol.~17, no.~1, pp. 53--69, Jan.\ 2008.

\bibitem{tropp2007omp}
J.~A.~Tropp and A.~C.~Gilbert,
``Signal recovery from random measurements via orthogonal matching
pursuit,''
\emph{IEEE Trans. Inf. Theory}, vol.~53, no.~12, pp. 4655--4666, Dec.\ 2007.

\bibitem{dong2011centralized}
W.~Dong, L.~Zhang, and G.~Shi,
``Centralized sparse representation for image restoration,''
in \emph{Proc. IEEE Int. Conf. Comput. Vis.\ (ICCV)}, 2011, pp. 1259--1266.

\bibitem{zhang2014groupbased}
J.~Zhang, D.~Zhao, and W.~Gao,
``Group-based sparse representation for image restoration,''
\emph{IEEE Trans. Image Process.}, vol.~23, no.~8, pp. 3336--3351, Aug.\ 2014.

\bibitem{tseng2001convergence}
P.~Tseng,
``Convergence of a block coordinate descent method for nondifferentiable
minimization,''
\emph{J. Optim. Theory Appl.}, vol.~109, no.~3, pp. 475--494, 2001.

\bibitem{bickel2009simultaneous}
P.~J.~Bickel, Y.~Ritov, and A.~B.~Tsybakov,
``Simultaneous analysis of Lasso and Dantzig selector,''
\emph{Ann. Statist.}, vol.~37, no.~4, pp. 1705--1732, 2009.

\bibitem{negahban2012unified}
S.~N.~Negahban, P.~Ravikumar, M.~J.~Wainwright, and B.~Yu,
``A unified framework for high-dimensional analysis of $M$-estimators
with decomposable regularizers,''
\emph{Statist. Sci.}, vol.~27, no.~4, pp. 538--557, 2012.

\bibitem{buhlmann2011statistics}
P.~B\"uhlmann and S.~van~de~Geer,
\emph{Statistics for High-Dimensional Data: Methods, Theory and
Applications.}\ Berlin: Springer, 2011.

\bibitem{donoho1994ideal}
D.~L.~Donoho and I.~M.~Johnstone,
``Ideal spatial adaptation by wavelet shrinkage,''
\emph{Biometrika}, vol.~81, no.~3, pp. 425--455, 1994.

\bibitem{beck2009fista}
A.~Beck and M.~Teboulle,
``A fast iterative shrinkage-thresholding algorithm for linear inverse
problems,''
\emph{SIAM J. Imaging Sci.}, vol.~2, no.~1, pp. 183--202, 2009.

\bibitem{chan2002mathematical}
T.~F.~Chan and J.~Shen,
``Mathematical models for local nontexture inpaintings,''
\emph{SIAM J. Appl. Math.}, vol.~62, no.~3, pp. 1019--1043, 2002.

\bibitem{bertalmio2000image}
M.~Bertalmio, G.~Sapiro, V.~Caselles, and C.~Ballester,
``Image inpainting,''
in \emph{Proc.\ ACM SIGGRAPH}, 2000, pp. 417--424.

\bibitem{jain2013low}
P.~Jain, P.~Netrapalli, and S.~Sanghavi,
``Low-rank matrix completion using alternating minimization,''
in \emph{Proc.\ ACM Symp.\ Theory Comput.\ (STOC)}, 2013, pp. 665--674.

\end{thebibliography}

\end{document}
